\section{NA WSTĘPIE OD TŁUMACZA}

*Kalevalę* zwykliśmy określać mianem epopei, eposu narodowego Finów. Nazwa ta, oznaczająca wielki poemat epicki, zakłada indywidualne lub też zbiorowe autorstwo, ustaloną wersję, zastosowany kanon kompozycyjny i stylistyczny. Epos kojarzy się nam z archaicznym, choć niekoniecznie prymitywnym tworzywem i jego potraktowaniem, z możliwie całościowym zasobem informacji, jaki zazwyczaj odlegli od nas w czasie i przestrzeni twórcy uznali za obraz godny przekazania pamięci pokoleń w postaci najpierw ustnej, a z czasem i zapisanej.

Dzisiaj, po dwustu latach rozwoju folklorystyki, zwielokrotniła się ilość dokonanych zapisów epiki, chociaż nie wszystkie mogą iść w zawody z zapisaną przed wiekami klasyką tego gatunku. Są jeszcze kraje, w których śpiewak czy opowiadacz uliczny gromadzi wokół siebie słuchaczy, pełniąc tę samą funkcję, co jego poprzednicy w zamierzchłej przeszłości. W kraju Kalevy będzie to jednak odtwórca tekstu pisanego. Ostatnia "opowiadająca z głowy" nosicielka tradycji, Marina Takalo, została odkryta i opisana wkrótce po zakończeniu II wojny światowej.

Jednakże w ciągu długiego życia Eliasa Lönnrota (1802-1884), człowieka, który powołał do istnienia Kalevalę, przewinęły się tysiące wyspecjalizowanych strażników pamięci plemiennej -- śpiewaków run i wierszowanych legend o świętych, płaczek pogrzebowych i wodzirejów weselnych, znawców zaklęć od chorób i wszelakich przypadłości w lesie, na lądzie i na wodzie, czyli po prostu znachorów i babek pomagających położnicom, dziadów odpustowych, wędrownych rzemieślników, sprzedawców ikon i księgarzy, wszelakiego rodzaju domokrążców, nie mówiąc już o zastępach nieprofesjonalnych użytkowników tego dziedzictwa. Była to wpisana w uświęcony tryb życia cząstka wiedzy, nie uchodzącej za taką w mieście czy na plebanii, ale niezbędnej na wszelki wypadek.

Przy całej różnorodności wyliczonych treści i zadań wykonawców obowiązywała jednolitość formy, tak zwane później metrum *Kalevali* -- *kalevalanmitta*, ośmiozgłoskowiec z dwoma, rzadziej trzema inicjalnymi akcentami, odmienny od rąbanego wiersza skaldów nordyckich, a także rozlewnej intonacji bylin i płaczów ruskich, zgodny natomiast ze strukturą fińszczyzny i wespół z dodatkowymi środkami mnemotechnicznymi jak aliteracja wewnątrz wersu i paralelizm treści wersów sąsiednich, organiczniej wiążący zdania i strofy niż na przykład rym w polskiej pieśni ludowej. *Kalevalanmitta*, dzięki zaadaptowaniu jej przez działające wśród ludu zakony żebracze -- franciszkanów i dominikanów -- do przekazywania pieśni religijnych i wierszowanych legend o świętych, a także modlitewek na różne okazje, tworzyła rodzaj wspólnego języka, pomostu między kulturami, a nawet nośnika synkretyzmu wierzeń religijnych. Był to sprawny mechanizm przekazu częstokroć archaicznych treści, nie do końca zrozumiałych dla przygodnych wykonawców. Przekłamania często zdarzają się w ustnych przekazach, zwłaszcza w wędrówce poprzez różne obszary dialektowe, jednakże sama ich więźba leksykalna, składniowa i rytmiczna była niezbyt podatna na nowinki dochodzące z obcojęzycznych ośrodków oddziaływania kulturalnego. Sam Elias Lönnrot w opublikowanym w roku 1840 zbiorze pieśni ludowych *Kanteletar* -- "Muza gęśli" albo, jak kto woli "Córka *kantele*" -- wyodrębnił niewielką garstkę importów słownych i melodycznych o charakterze przeważnie balladowym, dając wyraźną przewagę tradycyjnemu sposobowi śpiewania. Wyglądało to jak swoista linia okopów kultury rodzimej, obejmująca praktycznie wszystkie terytoria plemienne bałtyckofińskie, to jest Savo, Häme, Kainuu, nawet wyspy językowe fińskich drwali stale mieszkających w Szwecji, różne narzecza karelskie, aż do obrzeży Morza Białego, Ingrów, Wepsów, Estów i Liwów, a pośrednio też Saamów (Lapończyków), którzy utracili swój pierwotny język i przejęli jeden z dialektów starofińskich, poddając go dość gruntownej obróbce fonetycznej. Dlatego tak szeroki zasięg miały wyprawy Lönnrota w poszukiwaniu najwybitniejszych wówczas śpiewaków run.

Z ich to ust udawało mu się nieraz zapisać kilkusetwierszowe całości epickie. Martti Haavio (1899-1973) -- folklorysta, religioznawca i poeta, autor opublikowanej w Polsce *Mitologii fińskiej*, nadaje tym wariantom nazwę *epyllion*, małych eposów. W każdym z nich te same czyny przypisywano różnym bohaterom, łączono dowolnie wątki i osoby. Nie były to szczątki jakiejś przed wiekami rozbitej całości, raczej coś w rodzaju mozaiki, a przede wszystkim żywa materia opowieści, podlegająca ciągłym przemianom i reinterpretacjom. Można ją było tylko uchwycić w ruchu i zafiksować w kształcie, jaki się wydawał najbardziej godny uwagi i poszanowania: wymagało to zatem twórczego zaangażowania ze strony zbieracza i kompilatora całości.

Pierwsze zabytki pisane języka fińskiego -- teksty liturgiczne -- pochodzą z końca XIII wieku. Dalsze wzbogacenie języka przyniosła Reformacja. Wierni otrzymali Nowy, a później Stary Testament, pojawiły się pieśni nabożne, a także zapisy pieśni biesiadnych.

W wieku XVII pod wpływem rosnących zainteresowań przedchrześcijańskim dziedzictwem Skandynawów zaczęto spisywać również folklor miejscowy. Rymowany katalog bóstw Finów i Karelów jako przestrogę przed uprawianiem ich kultu podał już Mikael Agricola, współbojownik Marcina Lutra, w przedmowie do tłumaczenia *Psałterza*, wydanego drukiem w r. 1551. Pojawiły się też spekulacje lokalnych erudytów na temat fińskiego panteonu i próby porównania go z nomenklaturą grecką i rzymską, co zaspokajało wprawdzie ich ambicje, ale nie przybliżało istoty rzeczy. Najciekawszymi dokumentami pozostają rękopiśmienne *Obrzędy z głową, czyli zaślubiny zwierza*, zapisane w Viitasaari w początkach wieku XVIII oraz *Cantio ursina* Petrusa Bånga z roku 1675. zawierające opis stypy po niedźwiedziu (który spotkamy w runie XLVI).

Pod koniec wieku XVIII na uniwersytecie w Åbo (Turku) Henrik Gabriel Porthan, uczony o rozległych horyzontach, rozprawą *De poësi fennica* zainicjował badania folklorystyczne prowadzone na poziomie naukowym i w atmosferze przychylnego zainteresowania. Przez Europę wieją już wiatry Oświecenia, empiryzmu, racjonalizmu. Dokonywany jest wielki remanent całego świata widzialnego, próba jego uporządkowania, zamknięcia w ściślejszych i bardziej pojemnych niż dotąd formułach.

Rozproszkowane politycznie Niemcy przeżywają okres "burzy i naporu". Poddani chcą być obywatelami, Naród, lud, przestaje być utożsamiany z państwem. Protestanccy pietyści stwierdzają, że każdy ma prawo chwalić Boga w swoim języku ojczystym. Ludoznawstwo, językoznawstwo porównawcze zaczyna stawiać pierwsze kroki. Francuski mit o skorumpowanej cywilizacji i szlachetnych dzikusach ustępuje miejsca szacunkowi dla prostego człowieka i jego duchowego bogactwa. Najostrzejszej krytyce podlega przywilej urodzenia. Nagłemu przewartościowaniu towarzyszy zrozumiała przesada. Kopciuszek śpiesznie przywdziewa dworskie ciuchy. Obok autentycznych pieśni Słowian Południowych i rzetelnie gromadzonych przez Johanna Gottfrieda Herdera *Stimmen der Völker in Liedern* równie bezkrytyczny poklask zyskują *Pieśni Osjana* i *Rękopis Królodworski*, których apokryficzny charakter nieprędko rozszyfrowano. Świadczyło to przede wszystkim o rosnącym zapotrzebowaniu na literaturę poza kręgiem dotychczasowych jej odbiorców.

Nowinki te docierały i pod niebo Finlandii. "Romantycy z Turku" żywo reagowali na odkrycia i hasła Herdera, braci Humboldtów, wreszcie braci Grimm, pozwalające im docenić i coraz gorliwiej gromadzić skarby rodzimej literatury ustnej. Studiujący w Uppsali Carl Axel Gottlund, skupiając uwagę na postaci Väinämöinena, w latach 1818-1821 ogłosił drukiem dwa zbiory pieśni ludowych, a jeszcze wcześniej wyraził (po szwedzku) nadzieję, że z ich literackiej obróbki "mógłby się zrodzić nowy Homer, Osjan albo *Pieśń o Nibelungach*". Rzucona idea pobudziła kilku ówczesnych twórców do prób epickich, zwłaszcza wokół wątku Väinämöinena, reprezentowanego w źródłach najróżnorodniej i najobficiej.

Równocześnie nastąpił przełom w statusie politycznym Finów. Po kolejnej przegranej wojnie. której poligon stanowiła Finlandia, w roku 1809 król szwedzki odstąpił ją cesarzowi Wszechrosji. Wierne, a traktowane po macoszemu feudum uzyskało autonomię od nowego suwerena, który, nie nadwątlając układów stanowych, mógł szachować dotychczasową szwedzkojęzyczną elitę -- ludem. Jak na to zareagowała, świadczy powiedzenie przypisywane Adolfowi Ivarowi Arvidssonowi, szwedzkojęzycznemu publicyście, którego za radykalne poglądy władze carskie wydaliły z Finlandii: "Nie jesteśmy Szwedami. Rosjanami nie chcemy zostać, a więc bądźmy Finami." Wielu ze znanych ludzi pozmieniało wówczas nazwiska, ale jeszcze więcej spośród nich -- swój stosunek do lekceważonego dotąd języka.

Nazwiska nie zmienił Elias Lönnroth (tak brzmi zapis w metryce chrztu według ówczesnej szwedzkiej pisowni), syn wiejskiego krawca, urodzony 9 kwietnia 1802 w "wikarówce" Sammati, należącej do rozrzuconej wśród jezior gminy Haarijärvi. Ówże krawiec, Fredrik Juhana Lönnroth, miał żywy umysł i nie raz układał wiersze "po prostemu" na różne okoliczności. Słyszano też, jak je wyśpiewywał.

Nie odziedziczył jednak ojcowskiego fachu nawiedzony chłopak, którego czterolatkiem mama nauczyła liter. Jako sześciolatek wyróżnił się w nauce szwedzkiego pod kierunkiem starej kobiety z sąsiedztwa, dwunastolatka oddano do dwuletniej szkoły w Tammisaari, którą ukończył jako prymus w grudniu 1815, zadziwiając swoimi postępami w szwedczyźnie. Do wiosny 1818 uczył się w szkole katedralnej w Turku, pożegnawszy się z nią musiał zarabiać na życie igłą w okolicach rodzinnych, w marcu 1820 przyjęto go do liceum w Porvoo, ale już w kwietniu jął się pracy w Hämeenlinna jako pomocnik aptekarza. Przepraktykował u niego dwa lata. Dość rzucić okiem na mapę jego *Wanderjahre*, by zdać sobie sprawę z jego energii i ruchliwości. Jesienią 1822 przyjęty na uniwersytet w Turku, latem 1823 zdaje egzamin stypendialny, 1824 publikuje pierwszy przekład poezji na szwedzki, zarabia jako nauczyciel domowy i przy redagowaniu też kandydackich stypendystów. W roku 1827 ogłasza pracę magisterską pod tytułem *De Väinämöine priscorum fennorum numine* (O Väinämöinenie, bóstwie dawnych Finów), w 1828 wyrusza w pierwszą wędrówkę jako zbieracz pieśni i zaklęć (w najbliższych pięciu latach odbędzie jeszcze trzy podobne wyprawy). w 1831 zostaje sekretarzem świeżo założonego Fińskiego Towarzystwa Literackiego. W 1832 jako wykształcony lekarz z bogatym stażem życiowym i zawodowym ogłasza dysertację doktorską *Om Finnarnes magiska medicin* (O magicznym lecznictwie u Finów). To dopiero początek jego prac heraklesowych. Najdłuższy, prawie dwuletni ciąg wypraw odbędzie w latach 1841-1842. Dla nas znamienny pozostanie punkt startu, dochodzenie do awansu i wytyczony kierunek.

W latach 1829-1831 publikuje *Kantele* -- pierwszą, czteroczęściową wersję zgromadzonych przez siebie pieśni ludowych, by w latach 1840-1841 nadać jej kształt ostateczny pod wzmiankowanym już tytułem *Kanteletar*, wraz z obszernym wstępem metodycznym i niezbędnym słowniczkiem wyrażeń gwarowych.

W rękopisie pozostał ukończony przezeń jesienią 1833 zbiór pieśni o Väinämöinenie i Lemminkäinenie, czyli *Prakalevala*, licząca około 5000 wersów. Miał już wtedy za sobą *summę* humanistycznego wykształcenia, jakiego mógł mu dostarczyć miejscowy uniwersytet, przed sobą -- wzorzec *Iliady* (której jedną pieśń zdążył przełożyć z oryginału), wewnątrz zaś zakodowane trzeźwe widzenie rzeczy i wystarczającą dozę *sisu* -- fińskiej zawziętości i ambicji. Nie zadowoliła go pierwsza drukowana w 1835 redakcja *Kalevali* z 12000 wersów: w 1849 opublikował tę, którą dziś czytamy, prawie dwukrotnie obszerniejszą.

Znam jedynie "ze słyszenia" archiwalne bruliony Lönnrota. Jest tego mnóstwo -- gigantyczny pomnik mrówczej pracy i konsekwencji w dążeniu do założonego celu. Musiał dochować wierności nagromadzonemu tworzywu, przyjąć reguły gry, nawyki myślenia przekazicieli run, rozbujane a niestrzyżone słownictwo, które sam zapisywał i pierwszy interpretował -- z szacunkiem i pokorą. Uczeni fińscy zbadali szczegółowo wszystkie te geologiczne nawarstwienia, jakie się złożyły na *kokoonpano* -- zebranie w całość, montaż, kompozycję *Kalevali*. Rozczarowało to niektórych entuzjastów, zwłaszcza tych, co uwierzyli w istnienie jakiegoś mitycznego prawzoru, który należało jedynie odtworzyć... Jak pisze profesor Lauri Honko, raptem jedna trzecia tekstów jest identyczna z pierwotnymi zapisami, głębokim zmianom zostały bowiem poddane język, pisownia, metrum, tkanka fabularna.

Wiedział, czego chce. Chciał pokazać w pełnym blasku nie tylko epickie fabuły, ale ocalić skarbiec, jaki dał się przechować w języku i obrzędzie, wszystko, czemu groziło zlekceważenie i zapomnienie.

Stawką w grze była godność całego, traktowanego dotąd przedmiotowo narodu, jego niespodziewana nobilitacja. Finowie znali swoją inność i płacili za nią sporą cenę, jak również za próby przypodobania się ościennym wyroczniom w dziedzinie kultury. Nie wiedzieli jeszcze, że to, co poczytywano im za winę, może stanowić tytuł do chwały. Ten humanistyczny erudyta i praktyczny lekarz jak nikt przed nim zrozumiał rolę szamana-wizjonera, przewodnika, uzdrowiciela i wychowawcy plemienia -- i wszedł w nią bez reszty na całe życie.

Był absolutnym "pozytywistą środków" -- jako urzędowy lekarz, słownikarz, paremiograf, wykładowca uniwersytecki, rzecznik uprawnień fińszczyzny w użytku publicznym i nauce, wydawca, publicysta i organizator budzącego się życia społecznego i kulturalnego.

Lönnrot korzystał obficie nie tylko z kontaktów z najwybitniejszymi wówczas śpiewakami run, jak Juhana Kainulainen, Ontrei Malinen, Arhippa Perttunen i jego syn Miihkali, oraz wieloma innymi, ale również z pomocy rosnącego zastępu zbieraczy, w głównej mierze Sakari Topeliusa starszego, a także najsłynniejszego z nich, D.E.D. Europaeusa, któremu *Kalevala* zawdzięcza cały wątek tragicznych perypetii Kullervo. Zgromadzonego materiału nie usiłował naginać do klasycystycznych szablonów: wszystko, co tylko się dało, przybliżył i uczłowieczył. Zdać się mogło. że przerzucił łuk tęczy od najbardziej archaicznych wyobrażeń kosmogonii do realiów kultury ludowej, jaką można było oglądać na własne oczy w jego czasach. Zdegradowane bóstwa stają się pełnokrwistymi ludźmi, rzec by można, że nawet fińska odmiana współczesnych ruchów feministycznych znajduje swoje wzory w bohaterkach *Kalevali*. Epos nic na tym nie stracił...

Efekt rzeczywiście okazał się piorunujący. Pierwszy przekład całości na język szwedzki, pióra M.A. Castréna, powstał już w roku 1841, dokonany prozą na francuski przez Léonzona Le Duca w 1845 -- na cztery lata przed wyjściem w świat wersji ostatecznej dzieła. Towarzyszy temu entuzjastyczny wykład Jakuba Grimma w Królewskiej Akademii Pruskiej w Berlinie.

Dokonany za pośrednictwem wersji Le Duca rymowany przekład fragmentów *Kalevali* pióra Seweryny Duchińskiej, zamieszczony w "Bibliotece Warszawskiej" z roku 1869, odcięty jest od głównych komponentów poetyki oryginału nieporównanie bardziej, niż działo się to na przykład z polskimi próbami przekładowymi poezji ludów słowiańskich. Taki brak bezpośredniej informacji miał się utrzymać u nas niestety dłużej. Przeszkody były nie tylko językowe. Kryteria przekładu poetyckiego we wszystkich trzech zaborach po olśniewających próbach Mickiewicza, Słowackiego i Norwida z rzadkimi wyjątkami spadły do poziomu najbanalniejszej poprawności rymotwórczej. Nie wiedziano, jak sobie poradzić z ośmiozgłoskowcem, który przy dominującej wówczas sylabotonice rychło przeradzał się w nieznośną "katarynę". Brakło właściwych narzędzi. Uwagę na to zwrócił już w swoich *Szkicach literackich* (1901, a nawet wcześniej, w "Przewodniku Naukowym i Literackim" 1882) Józef Tretiak, szczerze zainteresowany popularyzacją fińskiego eposu.

Gdzie indziej natomiast można powiedzieć o kolejnej fali "burzy i naporu". To bezpośrednio pod oddziaływaniem *Kalevali* pojawiały się takie manifestacje epickie jak amerykańska *Pieśń o Hajawacie* (1855) Henry'ego Wadswortha Longfellowa, estoński *Kalevipoeg* (1862) Friedricha Reinholda Kreutzwalda, wreszcie łotewski *Lačplesis* (1888) Andreja Pumpura. Namnożyło się również przekładów, z których niejeden stał się wydarzeniem literackim w swoim kraju, jak niemiecki Antona Schieffnera (1852), czeski Josefa Holečka (1895), angielski W. F. Kirby'ego (1907). Powstanie wreszcie fińskiej szkoły folklorystycznej wytworzyło międzynarodową płaszczyznę naukową, w której obecność *Kalevali* stanowiła punkt odniesienia dla rosnącego zastępu badaczy antropologii kulturowej.

*Kalevala* funkcjonuje nie tylko w warstwie słownej. Nasze potoczne o niej wyobrażenie ukształtował swoją wszechstronną działalnością ilustratorską Akseli Gallen-Kallela (1865-1931). W wielu encyklopediach i wydawnictwach przeglądowych sztuki europejskiej rzucają się w oczy reprodukcje jego płócien, fresków, motywów graficznych, takich zwłaszcza jak: Kucie Sampo, Matka Lemminkäinena nad ciałem syna, Kullervo wygłaszający klątwę, Louhi wydzierająca Sampo ze statku Väinämöinena. Przemożna indywidualność artysty wyparła bez reszty poprzednie akademickie propozycje. Zamysł Lönnrota odczytywano plastycznie już od pierwszych wersji eposu. Była to nie tylko ewolucja nurtów i szkół artystycznych, ale całego widzenia historii, przełom w świadomości narodu, sprowadzenie dążeń i nadziei do wspólnego mianownika, odnalezienie stylu. Narzuca się porównanie z Wyspiańskim -- od kartonów *Iliady* do scenografii *Wesela* -- i zuchwałe przypuszczenie, że gdyby zaistniała recepcja *Kalevali* w Polsce z zachowaniem godziwych terminów, czyli około stu lat temu, niechybnie ilustrowałby ją właśnie on -- wizjoner, dla którego baśń i tradycja nie były nigdy ucieczką do raju dzieciństwa, lecz powrotem do rzeczywistości prawdziwszej niż oglądana w lustrze czy przez okno.

Rzecz jasna, że odrębność strukturalna fińszczyzny wobec rodziny języków indoeuropejskich nastręczała tłumaczom nieco więcej problemów niż dotychczasowa praktyka translatorska z tłumaczeniami klasyki starożytnej włącznie. Dodajmy do tego wymogi "ludowości", jak również pewnej "dziwności" w bezpośrednim odbiorze. Niemcy, Anglicy lub Czesi mogli dokonywać zapożyczeń z pomników literatury średniowiecznej, Rosjanie z zapisów bylin i bezpośrednio z gwary Pomorów, sąsiadujących o miedzę z Finami. U nas próby wywołania podobnego efektu doprowadzały do takich rezultatów, jak przekład *Iliady* na góralszczyznę. Francuzi, których do czasów Boya inteligencja czytała w oryginale, musieli się utrzymać w swoim akademickim kanonie absolutnej zrozumiałości, stąd korzystanie właśnie z pośrednictwa ich przekładów skazywało tłumacza na brak jakichkolwiek odniesień do bogactwa fińskich dialektów. Skandynawów tłumaczono u nas przez wersje niemieckie. Nic dziwnego, że nie znany był u nas niezawodny klucz do zasobu leksykalnego *Kalevali*, jaki stanowił opracowany przez samego Lönnrota monumentalny słownik fińsko-szwedzki *Suomalais-Ruotsalainen Sanakirja* (1866-1880). Byliśmy skazani na pośrednictwo "osób trzecich" ze wszystkimi fatalnymi konsekwencjami. Dotyczy to zarówno dziewiętnastowiecznych prób przekładowych Seweryny Duchińskiej i Franciszka Jezierskiego, jak i międzywojennych prób Jana Brzechwy, wyraźnie uzależnionych od słownictwa i emocjonalnej aury klasycznego przekładu na rosyjski L.P. Bielskiego (1889).

Z oryginału próbowała przekładać Maria Krahelska-Tołwińska, która w czasie pobytu w Finlandii opanowała bieżącą fińszczyznę (co jednak nie było równoznaczne ze znajomością języka eposu), oraz nie tak dawno zmarła Kazimiera Zawistowicz-Adamska, tłumacząca z przygotowaniem etnograficznym i niewątpliwą kompetencją. Wypada jedynie żałować, że poprzestała na tych niewielu próbach.

Poetycki przekład Józefa Ozgi-Michalskiego ma niewątpliwe uroki. Wskutek jednak braku bezpośrednich odniesień językowych jest on niczym ów mityczny olbrzym Anteusz oderwany od podłoża.

Klawiatura dźwiękowa fińszczyzny zda się nader odległa od naszej, niemniej w każdej technice muzycznej możliwe bywają transpozycje. To przeświadczenie ośmieliło mnie do podjęcia kolejnej próby. Nie moją rzeczą sądzić, czy mi się powiodła. Zawsze jest szansa dla kogoś, kto postara się to zrobić czytelniej, celniej i pewniejszą ręką. Tego, co ośmielam się zaprezentować dziś ewentualnemu czytelnikowi, nie uważam bynajmniej za rzecz skończoną. *Ars longa, vita brevis*. Dawno już wyszedłem z wieku studenta czy ochoczego rekruta, którym pozostawałem dość długo, jednakże podjąłem się tego niespodziewanego dla mnie zadania jako rodzaju rozpoznania z marszu i tylko w tej mierze zostało ono wykonane.

Ale i przy tej okazji winienem podziękować wszystkim, którzy mi w tym dopomogli. Najpierw Temu, który przydziela nam wszystkim zadania, chociaż nie zawsze je sobie wyobrażamy, nieraz od nich odbiegamy i nie do końca wykonujemy, od czasów praojca Adama powołując się na różne okoliczności łagodzące.

A później dobrym ludziom, którzy dawali mi odpowiedniego szturchańca: druhnie Teodorze Dąbrowskiej, której zawdzięczam pierwsze zetknięcie z Finlandią, profesorowi Pertti Virtaranta, który nie zraził się moim lekkomyślnym dyletantyzmem, tłumaczce Cecylii Lewandowskiej za lekcje uporu i wierności obranemu celowi. Tam, gdzie są dziś, trudno telefonować...

A na ostatek wszystkim doradcom łącznie z nieustępliwą Joanną Trzcińską-Mejor oraz moim rodzonym synem Łukaszem, który zaopatrzył mnie w komputer i wdrożył w jego używanie, a razem z nim -- nauczycielom i sponsorom, których miałem mnóstwo: bo Finowie mogą być dumni z wielu rzeczy, ale *Kalevali* też nie muszą się wstydzić i chętnie dzielą się nią z wszystkimi, którzy tego zapragną.